% META: {"id":"1NSI-TC-EVAL-A-corrige","chapitre":"1NSI-TYPES-CONSTRUITS","type_objet":"corrige_evaluation","evaluation_ref":"1NSI-TC-EVAL-A","capacites":["1NSI-TYPES-CONSTRUITS-C1","1NSI-TYPES-CONSTRUITS-C2","1NSI-TYPES-CONSTRUITS-C3","1NSI-TYPES-CONSTRUITS-C4","1NSI-TYPES-CONSTRUITS-C5"],"mode_creation":"ex_nihilo","sources_inspiration":[],"status":"approved","origin":"needs_review","status_history":["needs_review","audited","accepted","approved"],"release_acceptance":"RELEASE_OWNER_BATCH_ACCEPTANCE","acceptance_closure_digest":"sha256:9b3ccf9a81c5520fbb7e03b7b2d3e7bf2057a02834b6d908bf3be5f49f3b553f","acceptance_receipt_digest":"sha256:4fad86f0282e0fa4e0419cca1ad6379ae7af5a280c04a4a90880f90a1c044242"}
% BEGIN-VERIFY
% joueurs = [{"pseudo": "Ace", "score": 150}, {"pseudo": "Bob", "score": 200}, {"pseudo": "Cat", "score": 180}]
% meilleur = joueurs[0]
% for j in joueurs:
%     if j["score"] > meilleur["score"]:
%         meilleur = j
% assert meilleur["pseudo"] == "Bob" and meilleur["score"] == 200
% meilleur["score"] = 0
% assert joueurs[1]["score"] == 0
%
% def au_dessus(scores, seuil):
%     return [score for score in scores if score > seuil]
%
% def moyenne(scores):
%     if len(scores) == 0:
%         raise ValueError("le tableau des scores est vide")
%     return sum(scores) / len(scores)
%
% def classement(scores):
%     tries = sorted(scores, reverse=True)
%     return [(rang + 1, tries[rang]) for rang in range(len(tries))]
%
% def nb_sombres(image, seuil):
%     total = 0
%     for ligne in image:
%         for pixel in ligne:
%             if pixel <= seuil:
%                 total += 1
%     return total
%
% def chercher(contacts, nom):
%     if nom in contacts:
%         return contacts[nom]
%     return "Inconnu"
%
% def ajouter(contacts, nom, numero):
%     contacts[nom] = numero
%
% assert au_dessus([150, 200, 180, 210, 170], 175) == [200, 180, 210]
% assert moyenne([150, 200, 180, 210, 170]) == 182.0
% try:
%     moyenne([])
%     leve = False
% except ValueError:
%     leve = True
% assert leve
% assert classement([150, 200, 180]) == [(1, 200), (2, 180), (3, 150)]
% image = [[100, 200, 100], [200, 50, 200], [100, 200, 100]]
% assert len(image) == 3 and len(image[0]) == 3
% assert image[1][1] == 50
% assert nb_sombres(image, 100) == 5
% contacts = {"Ali": "55123456", "Sara": "55234567"}
% assert chercher(contacts, "Ali") == "55123456"
% assert chercher(contacts, "Zoe") == "Inconnu"
% ajouter(contacts, "Zoe", "55999999")
% assert contacts["Zoe"] == "55999999"
% carnet1 = {"Ali": "55123456"}
% carnet2 = carnet1
% carnet2["Sara"] = "55234567"
% assert carnet1 == {"Ali": "55123456", "Sara": "55234567"}
% END-VERIFY

\begin{corrige}{1NSI-TC-EVAL-A-EX1}
\textbf{Q1.} Le programme affiche \lstinline{Bob} puis \lstinline{200}.

En détail : \lstinline{meilleur} vaut d'abord le dictionnaire d'\lstinline{Ace}, de score $150$. Au premier tour, \lstinline{Ace} n'est pas strictement supérieur à lui-même. Au deuxième, $200 > 150$ : \lstinline{meilleur} devient \lstinline{Bob}. Au troisième, $180 > 200$ est faux, et \lstinline{Bob} reste retenu.

\textbf{Q2.} Le programme recherche le joueur ayant le meilleur score et affiche son pseudo puis son score.

\textbf{Q3.} \lstinline{joueurs[1]["score"]} vaut alors $0$.

La boucle n'a pas copié le dictionnaire de \lstinline{Bob} : elle a fait de \lstinline{meilleur} un second nom pour ce même objet, qui est aussi \lstinline{joueurs[1]}. Un dictionnaire étant mutable, la modification faite à travers l'un des deux noms est visible par l'autre.
\end{corrige}

\begin{corrige}{1NSI-TC-EVAL-A-EX2}
\textbf{Q1.}
\begin{python}
def au_dessus(scores, seuil):
    return [score for score in scores if score > seuil]
\end{python}
La comparaison est stricte, conformément à l'énoncé : un score égal au seuil n'est pas retenu.

\textbf{Q2.}
\begin{python}
def moyenne(scores):
    if len(scores) == 0:
        raise ValueError("le tableau des scores est vide")
    return sum(scores) / len(scores)
\end{python}
Le test de vacuité doit précéder le calcul : sans lui, la division renverrait une \lstinline{ZeroDivisionError}, moins explicite que l'erreur demandée.

\textbf{Q3.}
\begin{python}
def classement(scores):
    tries = sorted(scores, reverse=True)
    return [(rang + 1, tries[rang]) for rang in range(len(tries))]
\end{python}
\lstinline{sorted} construit une nouvelle liste et laisse \lstinline{scores} intact. Le rang affiché commence à $1$, alors que les indices commencent à $0$ : d'où le \lstinline{+ 1}.
\end{corrige}

\begin{corrige}{1NSI-TC-EVAL-A-EX3}
\textbf{Q1.} L'image comporte $3$ lignes et $3$ colonnes. Le nombre de lignes est \lstinline{len(image)} ; le nombre de colonnes est \lstinline{len(image[0])}.

\textbf{Q2.} Le pixel central vaut $50$. L'expression est \lstinline{image[1][1]} : le premier indice désigne la ligne, le second la colonne, et la numérotation commence à $0$.

\textbf{Q3.}
\begin{python}
def nb_sombres(image, seuil):
    total = 0
    for ligne in image:
        for pixel in ligne:
            if pixel <= seuil:
                total += 1
    return total
\end{python}
Deux boucles imbriquées sont nécessaires : la première parcourt les lignes, la seconde les pixels de chaque ligne. Sur l'exemple, \lstinline{nb_sombres(image, 100)} vaut $5$ : les quatre coins à $100$ et le centre à $50$.
\end{corrige}

\begin{corrige}{1NSI-TC-EVAL-A-EX4}
\textbf{Q1.}
\begin{python}
def chercher(contacts, nom):
    if nom in contacts:
        return contacts[nom]
    return "Inconnu"
\end{python}
Interroger directement \lstinline{contacts[nom]} lèverait une \lstinline{KeyError} pour un nom absent : le test d'appartenance est donc indispensable.

\textbf{Q2.}
\begin{python}
def ajouter(contacts, nom, numero):
    contacts[nom] = numero
\end{python}
Oui, la fonction modifie le dictionnaire reçu en argument. Un dictionnaire est mutable et n'est pas copié lors de l'appel : la fonction agit sur l'objet du programme appelant, ce qui explique qu'elle n'ait rien à renvoyer.

\textbf{Q3.} \lstinline{carnet1} vaut \lstinline|{"Ali": "55123456", "Sara": "55234567"}|.

L'affectation \lstinline{carnet2 = carnet1} ne crée aucune copie : elle donne un second nom au même dictionnaire. L'ajout effectué par \lstinline{carnet2} est donc visible par \lstinline{carnet1}. Pour obtenir deux carnets indépendants, il faudrait écrire \lstinline{carnet2 = dict(carnet1)}.
\end{corrige}
